\documentclass[12pt]{article}
\usepackage[document]{ragged2e}
\usepackage[letterpaper, margin=1in]{geometry}

\usepackage{amsmath}
\usepackage{amsthm}
\usepackage{amssymb}
\usepackage{booktabs} % fancy lines for tables
\usepackage{caption}
\usepackage{enumitem} % customize item labels
\usepackage{endnotes}
\usepackage{float}    % force figure position to match source
\usepackage{fontspec}
\usepackage{unicode-math}
\usepackage{graphicx}
\usepackage{listings}
\usepackage{pgfplots} % graph functions
\usepackage{setspace}
\usepackage{shellesc}
\usepackage{tabularx} % arbitrary table column widths
\usepackage{tikz}     % general drawing

\widowpenalty10000 % prevent widow and orphan lines
\clubpenalty10000

\defaultfontfeatures{
    Ligatures = TeX,
    Scale = MatchLowercase,
}

\setmainfont{LibertinusSerif}[
    Path            = fonts/Libertinus/,
    Extension       = .otf,
    UprightFont     = *-Regular,
    ItalicFont      = *-Italic,
    BoldFont        = *-Bold,
    BoldItalicFont  = *-BoldItalic]

\setmonofont{JetBrainsMono}[
    % Scale           = MatchLowercase,
    Path            = fonts/JetBrains_Mono/static/,
    UprightFont     = *-Regular,
    ItalicFont      = *-Italic,
    BoldFont        = *-Bold,
    BoldItalicFont  = *-BoldItalic]

\setmathfont{LibertinusMath-Regular}[
    Path            = fonts/Libertinus/,
    Extension       = .otf]

\pgfplotsset{compat=1.18, grid=both, axis lines=middle}
\lstset{basicstyle=\ttfamily,columns=fullflexible,breaklines=true,frame=shadowbox,keepspaces}

\newcommand{\nth}{^{\text{th}}}

\renewcommand{\qedsymbol}{$\blacksquare$}
\renewcommand{\implies}{\rightarrow}

\usetikzlibrary{arrows,automata,shapes}

\input{modules/dot.tex}
\input{modules/kicad.tex}
\input{modules/puml.tex}


\title{Assignment Name \\ \large Course Name \\ Section Number}
\author{Matthew McCall}

\begin{document}

\maketitle
\doublespacing{}

\section{Hello World!}
This template is a Lua\LaTeX-first article setup that ships with OpenType Libertinus text fonts and a full OpenType math font (via \lstinline|unicode-math|), so Lua\LaTeX can use modern math glyphs and features out of the box. It also includes \emph{batteries-included} diagram embedding for Graphviz DOT (\lstinline|\embeddot|), PlantUML (\lstinline|\embedpuml|), and KiCad schematics (\lstinline|\embedkicad|): small Lua helpers rebuild outputs only when sources change and cache the generated PDFs under \lstinline|out/| for fast reruns. To use all integrations, install \lstinline|lualatex| (a TeX distribution), \lstinline|latexmk|, and \lstinline|inkscape|; for specific generators you’ll also need \lstinline|dot| (Graphviz), \lstinline|plantuml| (and Java), and \lstinline|kicad-cli|. Shell-escape must be enabled for the external tool calls (see \lstinline|.latexmkrc|).

\section{Graph}
The document embeds externally generated diagrams via the Lua helpers; \lstinline|dot| produces the source SVG, Inkscape converts it to PDF, and the LaTeX macro simply renders the finalized file. The width percent you pass to \lstinline|\embeddot| is interpreted relative to \lstinline|\linewidth|, so the figure flexibly scales without touching the Lua scripts.

\begin{figure}[htbp]
  \centering
  \embeddot{graphs/example}
  \caption{Graphviz output from \lstinline|graphs/example.dot| produced through the shared Lua builder and displayed in the document.}
\end{figure}

The surrounding paragraph keeps the figure grounded and explains how the macro resolves the \lstinline|.dot| file stored beside the source.

\section{Embedding Pipeline}
The same shared builder drives each integration: it detects stale PDFs, runs the appropriate CLI (PlantUML for the workflow diagram, \lstinline|dot| for graphs, etc.), and returns the resulting path so the TeX macro can include it.

\begin{figure}[h]
  \centering
  \embedpuml[0.55]{graphs/embed_workflow}
  \caption{The Lua helper orchestrates CLI invocation and converts PlantUML output into a PDF ready for \LaTeX.}
\end{figure}

\section{Code Example}
The listing below showcases the JetBrains Mono font by compiling a minimal \lstinline|Hello, World!| C program. It is centered and captioned to keep the documentation consistent with the other figures.

\begin{singlespacing}
    \begin{lstlisting}[language=C, caption={Minimal C program demonstrating the embedded code listing font.}]
#include <stdio.h>
#include <stdlib.h>

int main(void) {
    printf("Hello, World!");
    return EXIT_SUCCESS;
}\end{lstlisting}
\end{singlespacing}

\end{document}
